\chapter{Introducción}
\label{cap:1}

\section{Contexto general del trabajo}

La gestión de emergencias es uno de los ámbitos de actuación pública donde la rapidez y la calidad de las decisiones tienen consecuencias más directas sobre la vida de las personas. En los primeros minutos de un incidente, el operador del 112 o el responsable de coordinación debe interpretar información incompleta, heterogénea y a menudo cambiante: texto libre de la llamada, categoría preliminar, localización, número de afectados, señales de riesgo, condiciones territoriales y posible necesidad de escalar la respuesta.

El marco español de protección civil, articulado principalmente por la Ley 17/2015 y la Norma Básica de Protección Civil, establece una respuesta basada en coordinación, planificación y actuación graduada ante diferentes riesgos \cite{Ley17_2015,RD524_2023}. En Castilla y León, este marco se concreta mediante la Ley 4/2007, el Plan Territorial de Protección Civil de Castilla y León (PLANCAL) y diversos planes especiales como INFOCAL, INUNCYL o MPCyL \cite{Ley4_2007_CYL,PLANCAL_DEC_4_2019,INFOCAL_DEC_6_2025,INUNCYL,MPCYL_ACUERDO_3_2008}. Este entorno normativo proporciona una base adecuada para construir un sistema de apoyo a la decisión que no se limite a clasificar incidentes, sino que justifique sus recomendaciones con criterios operativos y trazabilidad documental.

El presente TFM se sitúa en la intersección entre gestión de emergencias e inteligencia artificial aplicada. Su objetivo es diseñar, implementar y evaluar un prototipo software modular que priorice incidentes 112 de Castilla y León en una escala académica P1--P4, ofreciendo al operador una recomendación explicable, revisable y respaldada por reglas, evidencias internas y trazabilidad documental. La escala P1--P4 utilizada en el trabajo no corresponde a una prioridad oficial del servicio 112, sino a una construcción académica necesaria para entrenar, evaluar y explicar el prototipo. El sistema no sustituye al operador humano ni moviliza recursos de forma autónoma: actúa como herramienta de apoyo a la primera decisión.

\section{Planteamiento del problema}

La saturación informativa en emergencias aparece cuando el volumen, la velocidad o la heterogeneidad de la información supera la capacidad del operador para procesarla con seguridad. Este fenómeno es crítico en situaciones con múltiples incidentes simultáneos, avisos incompletos o información contradictoria. En esos escenarios, una priorización deficiente puede retrasar la atención a incidentes graves, sobredimensionar sucesos menores o dificultar la justificación posterior de la decisión adoptada.

El problema abordado por este TFM puede formularse así: dado un conjunto de incidentes 112 descritos mediante texto libre y datos operativos parciales, ¿cómo puede un sistema inteligente ayudar a estructurar la información, extraer variables relevantes, recomendar una prioridad P1--P4 y explicar esa recomendación de forma comprensible y auditable?

La respuesta propuesta se apoya en una arquitectura de tres capas. La Capa 1 transforma texto operativo en variables y señales predecisionales mediante procesamiento del lenguaje natural auditable, detección de negaciones y normalización operativa. La Capa 2 recomienda la prioridad mediante RuleFit-lite y se compara con una línea base heurística y una implementación alternativa de la misma familia de modelos. La Capa 3 genera una explicación mediante un modelo de lenguaje de gran tamaño (\textit{Large Language Model}, LLM) local, generación aumentada por recuperación (\textit{Retrieval-Augmented Generation}, RAG) sobre el corpus normativo y herramientas basadas en \textit{Model Context Protocol} (MCP). Esta separación impide que el componente generativo recalcule la prioridad o sustituya el criterio profesional.

\section{Justificación e impacto social}

Desde la perspectiva operativa, cualquier mejora en la primera decisión puede contribuir a orientar antes la atención hacia incidentes con riesgo vital, víctimas graves, atrapamiento, intoxicación, propagación o necesidad de coordinación compleja. El valor del sistema no reside en automatizar la emergencia, sino en reducir carga cognitiva y mejorar la consistencia de la recomendación inicial.

El \textit{Time To First Decision} (TTFD) se utiliza únicamente como motivación del problema. El prototipo no mide ni demuestra una reducción de ese tiempo en una sala 112 real o con operadores profesionales. La evaluación se centra en la trazabilidad, el rendimiento del modelo de priorización, la coherencia funcional del prototipo y la fidelidad de las explicaciones.

Desde la perspectiva tecnológica, las emergencias constituyen un dominio exigente para la IA: información incompleta, alto impacto potencial, necesidad de trazabilidad, supervisión humana obligatoria y prohibición de depender de servicios externos no controlados para datos sensibles. Esta orientación es coherente con los principios de transparencia, supervisión humana, gestión del riesgo y robustez recogidos en el Reglamento Europeo de Inteligencia Artificial \cite{EUAIAct2024}, así como con revisiones recientes sobre IA en gestión de emergencias \cite{SAPEA2025,Ghaffarian2023}.

Desde la perspectiva social, el proyecto se alinea con una visión responsable de la IA: una herramienta que amplifica la capacidad de análisis del operador, pero no desplaza su criterio profesional. En un dominio crítico, la explicabilidad no es decorativa; es una condición para que la recomendación pueda revisarse, discutirse, corregirse y auditarse.

\section{Contribución específica desde la inteligencia artificial}

La contribución principal del TFM no consiste en entrenar un modelo opaco para maximizar una métrica aislada de precisión. El proyecto aborda un problema más delicado: construir un sistema explicable y evaluable cuando no existe una etiqueta oficial P1--P4 en el conjunto de datos abierto. Por ello, la prioridad académica se construye mediante supervisión débil, anclada en PLANCAL y en la guía de etiquetado del proyecto.

La aportación de inteligencia artificial comprende la extracción auditable de variables desde texto, la construcción de etiquetas mediante supervisión débil, la priorización interpretable con RuleFit-lite, su contraste con una línea base heurística y con una alternativa externa de la familia RuleFit, y la generación de explicaciones a partir de la salida estructurada de Capa 2. Estas funciones se integran sin transferir al LLM la recomendación de prioridad ni la decisión final.

Esta combinación permite que cada recomendación pueda vincularse con la información de entrada, las variables extraídas, las reglas activadas, el grado de confianza, la trazabilidad documental y la decisión humana final. Esa trazabilidad constituye la diferencia entre un simple clasificador y un sistema de apoyo a la decisión.

\section{Finalidad y alcance del sistema propuesto}

El objetivo general del TFM es diseñar, implementar y evaluar un prototipo software modular de apoyo a la decisión para priorización temprana de incidentes 112 en Castilla y León. La finalidad es ofrecer una recomendación P1--P4 explicable y auditable, basada en variables operativas, aprendizaje interpretable, corpus normativo y supervisión humana.

\subsection{Alcance funcional}

El sistema abarca la recepción de un incidente, la extracción de variables, la recomendación de prioridad, la generación de una explicación operativa, el registro auditable de la inferencia y la posibilidad de que el operador acepte, modifique o rechace la recomendación. El alcance evaluado prioriza el núcleo científicamente defendible: variables operativas reproducibles, señales textuales, priorización interpretable, explicación normativa local y una interfaz de demostración para revisar el flujo completo. Adicionalmente, el prototipo incorpora una entrada de voz experimental que transcribe audio de prueba y propone campos estructurados para revisión; no se ha validado con llamadas reales.

La salida de Capa 2 incluye prioridad recomendada, probabilidades, reglas activadas y contexto explicativo. Capa 3 utiliza esa información para explicar la recomendación sin recalcularla, y el operador conserva la decisión final: puede aceptarla, modificarla o rechazarla.

No forman parte del alcance el despacho automático de recursos, la integración productiva con sistemas cerrados del 112, el uso de APIs cloud para inferencia con datos sensibles, la certificación regulatoria del sistema ni la validación externa con personal del 112. Estos elementos se documentan como trabajo futuro cuando corresponde.

\subsection{Alcance territorial y empírico}

El alcance territorial se centra en Castilla y León. El Registro de Emergencias del 112 de Castilla y León aporta la base empírica real para trabajar con incidentes, texto operativo, fechas, provincias, localidades y coordenadas \cite{Registro112CyL}. El marco normativo autonómico se apoya en Ley 4/2007, PLANCAL y planes especiales CyL. Esta decisión evita mantener una separación artificial entre marco normativo y conjunto de datos, y refuerza la coherencia del proyecto.

La escala P1--P4 mantiene en todo el trabajo su carácter académico y no equivale jurídicamente a las situaciones operativas oficiales.

\subsection{Limitaciones}

El TFM presenta varias limitaciones explícitas. El prototipo no se ha desplegado en una sala 112 real. El conjunto de datos abierto no contiene prioridad oficial P1--P4, por lo que la etiqueta se construye mediante supervisión débil y requiere análisis de acuerdo. Las variables V08--V11, relacionadas con meteorología, inundabilidad y riesgo tecnológico, quedan diferidas por coste de integración y por el objetivo de mantener reproducible el núcleo evaluado.

Estas limitaciones no invalidan el trabajo; delimitan su alcance. El valor del TFM reside en demostrar una arquitectura explicable, trazable y evaluable sobre datos reales abiertos, con un marco normativo claro y con supervisión humana como principio de diseño.

\section{Enfoque del TFM como desarrollo software}

El trabajo se desarrolla bajo la modalidad de TFM de Tipo 2, orientado al desarrollo de software. En coherencia con esta modalidad, el resultado no es solo una memoria teórica: se entrega un prototipo modular, un corpus normativo, pruebas reproducibles y evidencias de evaluación.

No obstante, la contribución académica no se reduce a programar una aplicación. El centro del trabajo es el diseño de un sistema inteligente responsable: cómo representar incidentes, cómo evitar fugas de información, cómo construir etiquetas académicas, cómo entrenar un modelo interpretable, cómo explicar sus recomendaciones y cómo evaluar si el resultado es defendible.

\section{Estructura de la memoria}

La memoria se organiza en diez capítulos y anexos técnicos. El Capítulo 1 introduce el problema, el alcance y la contribución. El Capítulo 2 revisa el contexto y el estado del arte en IA para emergencias, DSS, NLP, supervisión débil, RuleFit, LLM locales, RAG y MCP. El Capítulo 3 fija el marco operativo y normativo de Castilla y León. El Capítulo 4 formula objetivos, preguntas de investigación y metodología. El Capítulo 5 define requisitos. El Capítulo 6 presenta la arquitectura de tres capas. El Capítulo 7 desarrolla datos, variables, etiquetado y leakage. El Capítulo 8 documenta el prototipo. El Capítulo 9 presenta la evaluación. El Capítulo 10 sintetiza conclusiones y trabajo futuro.

Los anexos recogen la taxonomía, las variables, la guía de etiquetado, el protocolo de revisión, la línea base heurística, los contratos, el corpus RAG, las instrucciones del LLM, las fichas de modelo, las capturas y la trazabilidad extendida.
